\documentclass[11pt]{article}
\usepackage[a4paper,margin=1in]{geometry}
\usepackage{amsmath,amssymb,amsthm,mathtools}
\usepackage{array,booktabs}
\usepackage{enumitem}
\usepackage{xcolor}
\usepackage[colorlinks=true,linkcolor=black,citecolor=black,urlcolor=blue]{hyperref}
\usepackage[expansion=false]{microtype}

% ---- epistemic tags ----
\definecolor{cforced}{rgb}{0.0,0.42,0.0}
\definecolor{ccomp}{rgb}{0.0,0.0,0.66}
\definecolor{copen}{rgb}{0.62,0.0,0.0}
\newcommand{\forced}{\textcolor{cforced}{\textsf{[forced]}}}
\newcommand{\computed}{\textcolor{ccomp}{\textsf{[computed]}}}
\newcommand{\open}{\textcolor{copen}{\textsf{[open]}}}

% ---- notation ----
\newcommand{\Z}{\mathbb{Z}}
\newcommand{\Q}{\mathbb{Q}}
\newcommand{\R}{\mathbb{R}}
\newcommand{\Cx}{\mathbb{C}}
\newcommand{\Mah}{M}
\newcommand{\muS}{\mu_S}
\newcommand{\Zfive}{\Z/5\Z}
\DeclareMathOperator{\charge}{charge}
\DeclareMathOperator{\arglat}{arg}

\theoremstyle{plain}
\newtheorem{theorem}{Theorem}[section]
\newtheorem{lemma}[theorem]{Lemma}
\newtheorem{corollary}[theorem]{Corollary}
\newtheorem{proposition}[theorem]{Proposition}
\theoremstyle{definition}
\newtheorem{definition}[theorem]{Definition}
\theoremstyle{remark}
\newtheorem{remark}[theorem]{Remark}

\setlength{\parskip}{3pt}

\title{\textbf{The $\boldsymbol{\Zfive}$ Case of the No--Salem Dichotomy}\\[4pt]
\large A charge--coupled golden floor that confines the pentagon sector,\\
the Smyth bound that sharpens the general floor, and\\
why $\mu(5)=2$ stays \emph{computed}}
\author{Echo--Squirrel Research\\[2pt]
\small A companion to \emph{The Emission--Gap Theorem}, \emph{Lehmer's Box},\\
\small and the charge--measure stack (\texttt{testing-plausible-claims.md})}
\date{2026-06-30}

\begin{document}
\maketitle

\begin{abstract}
\noindent
The no--Salem dichotomy --- ``no charge--admissible algebraic integer has Mahler measure in the
band $(1,\muS)$, $\muS=1.3247\ldots$ the plastic number'' --- is the single lemma whose general
form would promote three currently \computed{} pillars of the stack (the universal floor
$\Mah\in\{1\}\cup[\varphi,\infty)$, the odd floor $\mu(\text{odd})=2$, and the empty odd band)
to \forced{}. The $\Z/3\Z$ case falls to an elementary argument because $2\cos120^\circ=-1$ is
rational. We ask the same of the first odd case where that technique breaks, $\Zfive$, where
$2\cos72^\circ=\varphi-1$ and $2\cos144^\circ=-\varphi$ are irrational and Galois--conjugate.

We obtain four \forced{} results and one sharpened \computed{} residual.
\textbf{(i)} The reciprocal/non--reciprocal split of \emph{Emission--Gap} Cor.~5.x, specialized:
every \emph{non--reciprocal} object of charge group exactly $\Zfive$ satisfies $\Mah\ge\muS$ by
Smyth, and every \emph{reciprocal} such object is $\Phi_5^{\,a}\times(\text{real reciprocal units})$
with $\Mah\in\{1\}\cup[\varphi^2,\infty)$. This \emph{improves the prior general forced floor from
the realification value $2^{1/5}=1.1487$ to $\muS=1.3247$}. \textbf{(ii)} The degree--four ``pure
pentagon'' (no real root) obeys the \emph{stronger} gap $\Mah\in\{1\}\cup[\varphi^4,\infty)$,
$\varphi^4\approx6.854$, via an explicit $\Q(\sqrt5)$ modulus coupling $t=\sigma(s)$. \textbf{(iii)}
The pure--power family $x^5-m$ realizes the floor: $\Mah=m\in\{1\}\cup[2,\infty)$, attaining
$\mu(5)=2$ at $x^5-2$. \textbf{(iv)} The full claim $\mu(5)=2$ --- i.e. $\Mah\notin(1,2)$ for
\emph{all} charge--$\Zfive$ objects --- stays \computed{}: the entire residual is now pinned to a
single set, the \emph{non--reciprocal} charge--$\Zfive$ objects in $[\muS,2)$, empty over the
verification window with boundary case $x^5-2$ at exactly $2$. Promotion is withheld under the
method's asymmetry rule; the residual is the reciprocal/Salem core of Lehmer's problem specialized
to the pentagon lattice, equivalently the totally--positive $(1,2)$ gap surviving fifth--root
pullback. Every non--asymptotic statement is re--derived in exact arithmetic.
\end{abstract}

\section{Introduction}

\subsection{The obligation this case is meant to discharge}
The charge--measure stack assigns to each monic $p\in\Z[x]$ a finite abelian \emph{charge group}
$\charge(p)$ (Definition~\ref{def:charge}) and studies the interaction of that invariant with the
Mahler measure $\Mah(p)$. Three results of the stack are currently asserted at \computed{} strength,
each a corollary of one universal statement:

\begin{quote}
\textbf{No--Salem dichotomy.} Every charge--admissible algebraic integer (one whose conjugate
arguments are all rational multiples of $\pi$) has $\Mah\notin(1,\muS)$; sharper, the realized
floor over charge--admissible objects is $\varphi$ for even charge and $2$ for odd charge.
\end{quote}

The dichotomy is exactly the place where Lehmer's problem is \emph{excised rather than solved}: a
Salem number has irrational on--circle arguments (it is charge--inadmissible by Kronecker--Niven),
so the charge restriction removes precisely the objects that populate $(1,\muS)$, and the floor that
results is the shadow of that removal, not a bound on Lehmer's infimum. This is the thesis of
\emph{Lehmer's Box}: confinement, not resolution. The present paper tests how much of the dichotomy
becomes \forced{} at $\Zfive$, the smallest odd modulus where the elementary $\Z/3\Z$ proof fails.

\subsection{Why $\Zfive$, and what changes}
A charge group $\Z/n\Z$ forces every conjugate angle into $(2\pi/n)\Z$. The $\Z/3\Z$ proof closes
because $2\cos120^\circ=-1\in\Z$: the complex pair's quadratic factor $x^2+\rho x+\rho^2$ is integral
the instant its modulus $\rho$ is, forcing $\rho\in\Z$ and collapsing the search to $x^3-2$ with
$\Mah=2$. \emph{Nothing couples.} At $\Zfive$ the cosines are irrational (\S\ref{sec:hard}), and
integrality can be repaired only by importing a Galois--partner angle class --- a genuine coupling
the odd case cannot avoid. The consequence, made precise below, is that the degree--four sector
gains a \emph{wider} forced gap than the dichotomy asks, while the floor--realizing sector retreats
one degree and re--exposes the same bedrock obligation.

\subsection{Result and discipline}
Table~\ref{tab:summary} states the outcome. The governing rule is the method's \emph{asymmetry rule}
(\texttt{testing-plausible-claims.md}, \S1): a finite search may \emph{demote} a universal claim or
\emph{extend} its status to \computed{} over a stated window, but may \emph{never promote} it to
\forced{}; only an attached structural proof promotes. Accordingly we report the forced sub--results
as theorems and hold the general claim at \computed{}, with the residual obligation named, not hidden.

\begin{table}[ht]
\centering
\small
\begin{tabular}{p{0.46\linewidth}p{0.28\linewidth}c}
\toprule
sub--class of charge--$\Zfive$ objects & measure & status\\
\midrule
non--reciprocal (general) & $\Mah\ge\muS=1.3247$ & \forced{} (Smyth)\\
reciprocal (general) & $\Mah\in\{1\}\cup[\varphi^2,\infty)$ & \forced{} real part\,/\,\computed{} window\\
degree--4 ``pure pentagon'' (no real root) & $\Mah\in\{1\}\cup[\varphi^4,\infty)$ & \forced{}\\
pure--power $x^5-m$ & $\Mah=m\in\{1\}\cup[2,\infty)$ & \forced{}\\
\textbf{all charge--$\Zfive$, all degrees} ($\mu(5)=2$) & $\Mah\notin(1,2)$ & \computed{} over window\\
\bottomrule
\end{tabular}
\caption{The $\Zfive$ ledger. The forced general lower bound is $\muS$ (improving the prior
realification value $2^{1/5}=1.1487$); the residual to $\mu(5)=2$ is the non--reciprocal class in
$[\muS,2)$.}
\label{tab:summary}
\end{table}

\section{Preliminaries}\label{sec:prelim}

\begin{definition}[Charge group; admissibility]\label{def:charge}
For squarefree monic $p\in\Z[x]$ with roots $\alpha_j=r_je^{2\pi i\theta_j}$, write each on--unit
argument $\theta_j\in\Q$ (when it is rational) as a reduced fraction with denominator $d_j$. The
\emph{charge group} is $\charge(p)=\Z/n\Z$ with $n=\operatorname{lcm}_j d_j$; if some $\theta_j$ is
irrational, $p$ is \emph{charge--inadmissible} ($\charge(p)=\bot$). $p$ is \emph{charge--admissible}
iff every conjugate argument is a rational multiple of $2\pi$.
\end{definition}

\begin{lemma}[Odd modulus forbids negative reals]\label{lem:noneg}
If $\charge(p)=\Zfive$ then $p$ has no negative real root, and its on--unit roots are fifth roots of
unity. \forced{}
\end{lemma}
\begin{proof}
The lattice $(2\pi/5)\Z=\{0^\circ,72^\circ,144^\circ,216^\circ,288^\circ\}$ omits $180^\circ$, so no
root has argument $\pi$: all real roots are positive. An on--unit root has modulus $1$ and argument
in the lattice, hence equals $e^{2\pi ik/5}$, a fifth root of unity, contributing a $\Phi_5$ factor.
\end{proof}

\begin{definition}[Mahler measure]
For $p=a_d\prod_j(x-\alpha_j)\in\Z[x]$, $\Mah(p)=|a_d|\prod_j\max(1,|\alpha_j|)$. It is multiplicative,
$\Mah(pq)=\Mah(p)\Mah(q)$, and equals $1$ exactly on products of cyclotomics and powers of $x$
(Kronecker).
\end{definition}

\begin{lemma}[Smyth's bound]\label{lem:smyth}
Every non--reciprocal monic $p\in\Z[x]$ with $\Mah(p)>1$ satisfies $\Mah(p)\ge\muS=1.3247179\ldots$,
the real root of $x^3-x-1$ (the plastic number, the smallest Pisot number). \forced{} \emph{(Smyth
1971; this is exactly Thm.~\textup{smythfloor} of the substrate paper, with the caveat of
Rem.~\textup{reciprocal} that the bound is unconditional only off the reciprocal locus.)}
\end{lemma}

A polynomial is \emph{reciprocal} if its root multiset is closed under $\alpha\mapsto\alpha^{-1}$,
equivalently $x^{\deg p}p(1/x)=p(x)$. Salem minimal polynomials are reciprocal; Lehmer's number is
reciprocal. The entire difficulty of Lehmer's problem, and of the dichotomy, lives on the reciprocal
side, because Smyth disposes of the other side outright.

\begin{lemma}[Real reciprocal units]\label{lem:realrecip}
A reciprocal pair of real algebraic conjugates $\{r,r^{-1}\}$ with $r>1$ has integer trace
$r+r^{-1}\ge3$ (the value $2$ forces $r=1$), hence $r\ge(3+\sqrt5)/2=\varphi^2$ and Mahler
contribution $\ge\varphi^2\approx2.618$. A totally real reciprocal unit therefore has
$\Mah\in\{1\}\cup[\varphi^2,\infty)$. \forced{}
\end{lemma}

\section{Why $\Zfive$ is the hard analogue of $\Z/3\Z$}\label{sec:hard}

\begin{proposition}[Irrational, Galois--coupled cosines]\label{prop:cos}
$2\cos72^\circ=\varphi-1$ and $2\cos144^\circ=-\varphi$, where $\varphi=(1+\sqrt5)/2$ and
$\varphi^2-\varphi-1=0$. Both are irrational, and they are the two roots of $x^2+x-1$, hence
Galois--conjugate over $\Q$ under $\sqrt5\mapsto-\sqrt5$. \forced{} \emph{(verified symbolically,
\S\ref{sec:verify})}
\end{proposition}

A single complex pair at $\pm72^\circ$ with rational modulus $s$ contributes the quadratic
$x^2-s(\varphi-1)x+s^2$, whose linear coefficient is irrational. Integrality cannot be achieved
within that pair; it can be repaired only by adjoining the $\pm144^\circ$ pair as the Galois partner,
so that the $\sqrt5$--parts cancel between the two factors. \emph{This forced coupling is the
structural feature $\Z/3\Z$ never needs}, and it is the source both of the wider degree--four gap
(\S\ref{sec:pentagon}) and of the obstruction in the general case (\S\ref{sec:residual}).

\section{The reciprocal/non--reciprocal split: the forced floor}\label{sec:split}

We specialize the floor argument of \emph{Emission--Gap} Cor.~5.x to the charge--$\Zfive$ lattice.
It is the cleanest forced statement available for the general case, and it strictly improves the
realification bound used in the prior pass.

\begin{theorem}[Forced general floor at $\Zfive$]\label{thm:floor}
Let $O$ be charge--admissible with $\charge(O)=\Zfive$ and $\Mah(O)>1$.
\begin{enumerate}[label=\textup{(\alph*)},leftmargin=2.2em,topsep=2pt]
\item If $O$ is non--reciprocal, then $\Mah(O)\ge\muS=1.3247179\ldots$ \forced{}
\item If $O$ is reciprocal, then $O=\Phi_5^{\,a}\cdot U$ with $U$ a product of reciprocal units off
the unit circle; its real part obeys $\Mah\ge\varphi^2$, and no charge--$\Zfive$ reciprocal object in
the verification window has $\Mah\in(1,\varphi^2)$. \forced{} (real part) / \computed{} (window)
\end{enumerate}
Consequently the forced lower bound for the non--reciprocal class --- which contains every realized
charge--$\Zfive$ object of measure below $\varphi^2$, including the floor object $x^5-2$ --- is
$\muS$, improving the realification value $2^{1/5}=1.1487$.
\end{theorem}

\begin{proof}
(a) Immediate from Lemma~\ref{lem:smyth}. (b) By Lemma~\ref{lem:noneg} the on--unit roots of $O$ are
fifth roots of unity, contributing the cyclotomic factor $\Phi_5^{\,a}$; strip it. The remaining
factor $U$ is reciprocal, charge--admissible, non--cyclotomic, and --- being charge--admissible ---
not Salem, so it has \emph{no} on--circle conjugate at irrational angle, and by the previous sentence
no on--circle conjugate at all. Thus every root of $U$ lies off the unit circle and is paired with
its inverse. The real such pairs are real reciprocal units, $\Mah\ge\varphi^2$ by
Lemma~\ref{lem:realrecip}. The complex such pairs sit at the lattice angles $\pm72^\circ,\pm144^\circ$
and are coupled in $\Q(\sqrt5)$ exactly as in \S\ref{sec:pentagon}; the exact--arithmetic window
(\S\ref{sec:verify}) exhibits no reciprocal charge--$\Zfive$ object with $\Mah<\varphi^2$, all
observed instances being $\Phi_5^{\,a}\times$(real reciprocal units) with $\Mah\in\{1,\varphi^2,
2+\sqrt3,\dots\}$. The realized minimum over the non--reciprocal class is $2$ (\S\ref{sec:purepower}),
which lies above $\muS$, so the forced bound $\muS$ is not attained and a gap remains
(\S\ref{sec:residual}).
\end{proof}

\begin{remark}[This is the corpus's own discipline]
Theorem~\ref{thm:floor} is the $\Zfive$ instance of the same case split that \emph{Emission--Gap}
Cor.~5.x runs for the spectral emission algebra (non--reciprocal $\to$ Smyth $\muS$; reciprocal
non--Salem $\to$ real reciprocal units $\to\ge\varphi$; realized at $\varphi$) and that the substrate
paper states as Thm.~\textup{smythfloor} with Rem.~\textup{reciprocal}. The reciprocal locus is, in
the words of that remark, ``the unsolved core''; here it is what keeps part~(b) at \computed{} rather
than fully \forced{}, and what blocks promotion of the general claim in \S\ref{sec:residual}.
\end{remark}

\section{The pure--pentagon theorem (degree four)}\label{sec:pentagon}

The degree--four sector, having no real root, is closed completely and with a \emph{wider} gap than
the dichotomy asks.

\begin{theorem}[Pure pentagon]\label{thm:pentagon}
Every irreducible quartic with charge group $\Zfive$ has $\Mah\in\{1\}\cup[\varphi^4,\infty)$,
$\varphi^4\approx6.854$; in particular $\Mah\notin(1,2)$, with the empty interval $(1,\varphi^4)$.
\forced{}
\end{theorem}

\begin{proof}
By Lemma~\ref{lem:noneg} and the no--real--root hypothesis, the object is a product of a $\pm72^\circ$
conjugate pair (common modulus $s$) and a $\pm144^\circ$ pair (common modulus $t$):
\[
O(x)=\bigl(x^2-s(\varphi-1)x+s^2\bigr)\bigl(x^2+t\varphi\, x+t^2\bigr),\qquad s,t>0.
\]
Expanding and using $\varphi^2-\varphi=1$ (Prop.~\ref{prop:cos}) to collapse the cross term,
\[
[x^3]=\varphi(t-s)+s,\quad [x^2]=s^2+t^2-st,\quad [x^1]=st\bigl(\varphi(s-t)+t\bigr),\quad [x^0]=(st)^2.
\]
Integrality of all four coefficients forces the $\sqrt5$--parts carried by $\varphi$ to cancel, which
holds iff $s$ and $t$ are Galois--conjugate in $\Q(\sqrt5)$: $t=\sigma(s)$ under $\sqrt5\mapsto-\sqrt5$.
Hence $s,t$ are the two positive roots of a totally real integer quadratic
\[
y^2-ky+m,\qquad k=s+t\in\Z\ (\text{trace}),\quad m=st\in\Z_{>0}\ (\text{norm}).
\]
With two roots at modulus $s$ and two at $t$, $\Mah=\max(1,s)^2\max(1,t)^2$. The case analysis on
$(k,m)$:
\begin{center}\small
\begin{tabular}{lccc}
\toprule
regime & constraint & measure & regime minimum\\
\midrule
$s,t>1$ & $m=st\ge2$ & $\Mah=m^2$ & $16$ at $s{=}t{=}2$, $(k,m){=}(4,4)$\\
$s>1>t,\ m=1$ ($t=1/s$) & $s=\tfrac{k+\sqrt{k^2-4}}2,\ k\ge3$ & $\Mah=s^2$ & $\boldsymbol{\varphi^4}$ at $k{=}3$ ($s{=}\varphi^2$)\\
$s>1>t,\ m\ge2$ & $s>m$ & $\Mah=s^2>m^2$ & $>\varphi^4$\\
$s,t<1$ & $st=m\ge1$ impossible & --- & excluded\\
$t=1$ or $s=1$ & on--circle root & reduces to $\Phi_5$ & $\Mah=1$\\
\bottomrule
\end{tabular}
\end{center}
The global minimum over the nontrivial regimes is $\varphi^4$, at $(k,m)=(3,1)$: $s=\varphi^2\approx
2.618$ at $\pm72^\circ$, $t=\varphi^{-2}\approx0.382$ at $\pm144^\circ$, giving the explicit minimizer
\[
x^4-x^3+6x^2+4x+1,\qquad \charge=\Zfive,\quad \Mah=\varphi^4.
\]
Several modulus pairs per angle factor the object into such Galois quartets and integer powers, across
which $\Mah$ multiplies, preserving membership in $\{1\}\cup[\varphi^4,\infty)$.
\end{proof}

\begin{remark}
The contrast with $\Z/3\Z$ is now exact: $\Z/3\Z$ couples nothing and bottoms out at $\Mah=2$;
$\Zfive$ couples the two moduli as $\Q(\sqrt5)$--conjugates and bottoms out \emph{higher}, at
$\varphi^4$. The coupling that makes the odd case hard in general is the very thing that widens the
degree--four gap.
\end{remark}

\section{The pure--power family and the realized floor}\label{sec:purepower}

\begin{theorem}[Pure power]\label{thm:purepower}
For $m\in\Z_{\ge2}$, $x^5-m$ has charge group $\Zfive$, is non--reciprocal, and has $\Mah=m$. Hence
the family $x^5-m$ realizes $\Mah\in\{1\}\cup[2,\infty)$, attaining the dichotomy floor $\mu(5)=2$ at
$x^5-2$. \forced{}
\end{theorem}
\begin{proof}
The roots are $m^{1/5}e^{2\pi ik/5}$, $k=0,\dots,4$: one positive real and two conjugate pairs at
$\pm72^\circ,\pm144^\circ$, all of modulus $m^{1/5}>1$, so $\charge=\Zfive$ and
$\Mah=(m^{1/5})^5=m$. Closure under $\alpha\mapsto\alpha^{-1}$ would require $m^{-1/5}$ among the
roots; it is not, so the family is non--reciprocal (consistent with Theorem~\ref{thm:floor}(a):
$m\ge2>\muS$).
\end{proof}

This is where the real root enters and the pentagon theorem stops applying: in $x^5-2$ the real root
$2^{1/5}$ is Galois--conjugate to the four complex roots and cannot be split off, so the object is
genuinely quintic and carries a real direction the degree--four analysis excluded.

\section{The computed residual and the obligation}\label{sec:residual}

\begin{proposition}[Location of the residual]\label{prop:residual}
Combining Theorems~\ref{thm:floor}--\ref{thm:purepower}: the reciprocal charge--$\Zfive$ class lies in
$\{1\}\cup[\varphi^2,\infty)$ (above $2$); the non--reciprocal class is forced into
$\{1\}\cup[\muS,\infty)$ with realized minimum $2$. Therefore the only way a charge--$\Zfive$ object
could land in $(1,2)$ is to be non--reciprocal with $\Mah\in[\muS,2)$. Over the verification window
(quartics $|c|\le10$, quintics $|c|\le4$, sextics $|c|\le3$) this set is empty; the boundary object
$x^5-2$ sits at exactly $2$. \computed{}
\end{proposition}

\begin{remark}[Why the residual does not promote]\label{rem:notforced}
Three handles are available and none reaches $2$.
\emph{(1) Realification.} $\psi^5(O)$ (fifth powers of the roots) is totally positive with
$\Mah(\psi^5O)=\Mah(O)^5$, so the totally--positive $(1,2)$ gap yields only $\Mah(O)\ge2^{1/5}=1.1487$
--- the fifth power is unavoidable, and asking $\psi^5O$ to clear $32$ instead of $2$ is the original
claim restated, hence circular.
\emph{(2) Smyth.} Lemma~\ref{lem:smyth} forces $\muS=1.3247$ on the non--reciprocal class --- the
best off--the--shelf bound, and a strict improvement on (1), but still below $2$.
\emph{(3) Classification.} The pentagon theorem closes degree four and the pure--power theorem closes
$x^5-m$, but a general non--reciprocal charge--$\Zfive$ object of degree $\ge5$ with non--constant
modulus structure is not classified.
The residual --- ``no non--reciprocal charge--$\Zfive$ object has $\Mah\in[\muS,2)$'' --- is the
reciprocal/Salem core of Lehmer's problem specialized to the pentagon lattice (substrate
Rem.~\textup{reciprocal}: no proven uniform positive floor on that locus), equivalently the
totally--positive $(1,2)$ gap surviving fifth--root pullback \emph{with} the rotation--class coupling
that $\psi^5$ destroys. That gap \emph{reduces into} the Schur--Siegel--Smyth (SSS) trace problem (Siegel~1945;
Smyth~1984), a recognized open problem: SSS--class progress on the totally--positive spectrum
would close the residual. Whether the pentagon residual is \emph{equivalent} to SSS or only
\emph{embeds into} it --- \textbf{the exact relationship (one--way reduction vs.\ equivalence)}
--- is itself unsettled, and we record it as an explicit \open{} sub--question rather than as
undifferentiated ``territory''. It is not closable from scratch in--session. By the asymmetry rule, a clean window plus the four forced sub--results is
the honest maximum; promotion of $\mu(5)=2$ to \forced{} still requires the trace--problem lemma,
attached and referenced.
\end{remark}

\section{Verification (exact arithmetic, \texorpdfstring{$\textsf{dps}=45$}{dps=45})}\label{sec:verify}

Every decision is symbolic or exact; floats only screen. Engine: the method's own
\texttt{charge\_group}, multiplicity--correct exact Mahler (integer factorization), and the
\texttt{np\_admis} angle pre--filter.

\begin{center}\small
\begin{tabular}{ll}
\toprule
check & result\\
\midrule
$2\cos72^\circ=\varphi-1$, \ $2\cos144^\circ=-\varphi$ & \textbf{True} (symbolic)\\
$\varphi^2-\varphi-1=0$ (cross--term collapse) & \textbf{0}\\
$\{2\cos72^\circ,2\cos144^\circ\}$ are roots of & $x^2+x-1$ $\Rightarrow$ Galois--conjugate over $\Q$\\
charge--$5$ quartics, $|c|\le10$: distinct $\Mah$ & $\{1,\varphi^4\}$, \ $\varphi^4=6.854102$\\
$(k,m)$ minimizer; minimizer object & $(3,1)\to s=\varphi^2$; \ $x^4-x^3+6x^2+4x+1$, $\Mah=\varphi^4$\\
charge--$5$ quartics with $\Mah\in(1,2)$ & \textbf{0}\\
non--reciprocal charge--$5$ found; $\min\Mah$ & $13$ objects; \ $\min=2$ ($x^5-2$ and kin)\\
all non--reciprocal obey $\Mah\ge\muS$ & \textbf{True}\\
non--reciprocal charge--$5$ in $[\muS,2)$ & \textbf{0} (the residual window)\\
reciprocal charge--$5$ found; measures & $5$ objects; \ $\{1,\varphi^2,2+\sqrt3\}$ (all $=\Phi_5\times$ real recip.\ unit)\\
reciprocal charge--$5$ in $(1,2)$ & \textbf{0}\\
$\psi^5(x^5-2)$ & $2$ (mult.\ $5$), totally positive; $\Mah(O)=32^{1/5}=2$\\
\textbf{overall charge--$5$ with $\Mah\in(1,2)$} & \textbf{0} \quad (realized floor $2$)\\
\bottomrule
\end{tabular}
\end{center}

\section{Verdict record and epistemic ledger}

\begin{small}
\begin{verbatim}
{
  "claim_id": "coupling/no-salem-dichotomy@Z5/v2-smyth",
  "statement": "every charge-admissible object with charge group exactly Z/5Z has
                M not in (1,2); equivalently mu(5)=2",
  "verdict": "forced [non-recip: M>=mu_S] · forced [deg-4 pentagon: M in {1}cup[phi^4,inf)]
              · forced [x^5-m: M=m>=2] · computed [general mu(5)=2]",
  "forced": {
    "general_lower_bound": "mu_S=1.3247179 (Smyth, non-reciprocal class);
                            improves prior realification 2^(1/5)=1.1487",
    "reciprocal_class": "Phi_5 x real-reciprocal-units; real part M>=phi^2 forced",
    "deg4_pure_pentagon": "M in {1} cup [phi^4, infty); Galois coupling t=sigma(s),
                           s,t roots of y^2-ky+m; min phi^4 at (k,m)=(3,1)",
    "pure_power": "x^5-m gives M=m in {1} cup [2, infty); mu(5)=2 at x^5-2"
  },
  "computed": {
    "realized_floor": 2,
    "residual": "non-reciprocal charge-Z/5Z in [mu_S, 2), empty over window;
                 boundary x^5-2 at exactly 2",
    "window": "quartics |c|<=10, quintics |c|<=4, sextics |c|<=3"
  },
  "proof_obligation": "non-reciprocal charge-Z/5Z in [mu_S,2) = reciprocal/Salem core
                       specialized; equiv totally-positive (1,2) gap surviving 5th-root
                       pullback with the rotation-class coupling psi^5 collapses",
  "engine": {"polynomial": "ok", "dps": 45, "angle_tol": "1e-18"},
  "corpus_refs": ["emission_gap cor:floor", "vector_substrate thm:smythfloor + rem:reciprocal",
                  "lehmers_box (sidestep)", "lehmersproblemanintroduction (Smyth/Dobrowolski)"],
  "provenance": "sha256:e1781b5f0dffd89132ed306fa8d09e722ccfe8a682ee5e365f86bcdb3be45b62"
}
\end{verbatim}
\end{small}

\begin{center}\small
\begin{tabular}{p{0.56\linewidth}cl}
\toprule
claim & status & backing\\
\midrule
Odd modulus forbids negative reals; on--unit $\Rightarrow$ fifth root of unity & \forced{} & Lem.~\ref{lem:noneg}\\
$2\cos72^\circ,2\cos144^\circ$ irrational, roots of $x^2+x-1$ & \forced{} & Prop.~\ref{prop:cos}, \S\ref{sec:verify}\\
Non--reciprocal charge--$\Zfive$ $\Rightarrow\Mah\ge\muS$ & \forced{} & Thm.~\ref{thm:floor}(a), Smyth\\
Reciprocal charge--$\Zfive$ real part $\Rightarrow\Mah\ge\varphi^2$ & \forced{} & Thm.~\ref{thm:floor}(b), Lem.~\ref{lem:realrecip}\\
\quad reciprocal complex part has no sub--$\varphi^2$ instance & \computed{} & window \S\ref{sec:verify}\\
Degree--4 pure pentagon $\Rightarrow\Mah\in\{1\}\cup[\varphi^4,\infty)$ & \forced{} & Thm.~\ref{thm:pentagon}\\
$x^5-m\Rightarrow\Mah=m$; $\mu(5)=2$ at $x^5-2$ & \forced{} & Thm.~\ref{thm:purepower}\\
Forced general floor improves $2^{1/5}\to\muS$ & \forced{} & Thm.~\ref{thm:floor}\\
$\mu(5)=2$ (no charge--$\Zfive$ object in $(1,2)$) & \computed{} & Prop.~\ref{prop:residual}, window\\
Residual embeds into SSS; exact relation (reduction vs.\ equivalence) & \open{} & Rem.~\ref{rem:notforced}\\
\bottomrule
\end{tabular}
\end{center}

\section{Position in the corpus}

This paper is a leaf of the charge--measure stack and inherits its vocabulary and discipline.
\begin{itemize}[leftmargin=1.4em,topsep=2pt,itemsep=1pt]
\item \emph{The Emission--Gap Theorem} --- the floor argument Cor.~5.x (non--reciprocal $\to$ Smyth,
reciprocal non--Salem $\to$ real reciprocal units, realized at $\varphi$) is specialized here in
Theorem~\ref{thm:floor}; Thm.~\textup{main} (no Salem emitted) is the closure this dichotomy makes
universal--in--charge.
\item \emph{Lehmer's Box} --- the sidestep logic: the box theorem is unconditional, strictly weaker
than Lehmer's conjecture, and consistent with both its truth values. The $\Zfive$ residual is the
same open problem entered through the pentagon door; we confine, we do not resolve.
\item \emph{The Occupant of the Salem Slot} --- the trace redirection $t=\theta+\theta^{-1}$ and the
$\sqrt5$ limit at the floor; the realification $\psi^5$ of \S\ref{sec:residual} is the charge--$5$
analogue of that boundary map, and loses the same coupling.
\item \emph{The substrate paper} (\texttt{vector\_substrate}) --- Thm.~\textup{smythfloor} is our
Lemma~\ref{lem:smyth}; Rem.~\textup{reciprocal} is the precise statement of why
Theorem~\ref{thm:floor}(b) and Proposition~\ref{prop:residual} stop at \computed{}.
\item \emph{The Operator Algebra of the Emission Semiring} --- the $\lambda$--ring with the Mahler
measure and the angle charge as its two characters; the charge group of this paper is the same
character, read at modulus $5$.
\item \emph{Lehmer's Problem: An Introduction} and the Method (\texttt{testing-plausible-claims.md})
--- Smyth/Dobrowolski background and the asymmetry rule that governs the forced/computed line drawn
throughout.
\end{itemize}

\paragraph{Closing.} The dichotomy did not fall at $\Zfive$; it bent. The pentagon sector gained a
\emph{wider} forced gap ($\varphi^4$) than asked, the general forced floor rose from $2^{1/5}$ to
$\muS$, and the residual contracted to a single named set --- the non--reciprocal charge--$\Zfive$
objects in $[\muS,2)$ --- which is the reciprocal/Salem core of Lehmer's problem wearing a pentagon.
$\mu(5)=2$ remains \computed{}, and that is the correct call: the window is clean, the sub--results
are forced, and the one statement that would promote the whole is exactly the trace--problem lemma we
do not have. Naming it is the result.

\begin{thebibliography}{9}
\small
\bibitem{Lehmer} D.~H. Lehmer, \emph{Factorization of certain cyclotomic functions}, Ann. of Math.
(2) \textbf{34} (1933), 461--479.
\bibitem{Kronecker} L.~Kronecker, \emph{Zwei S\"atze \"uber Gleichungen mit ganzzahligen
Coefficienten}, J. Reine Angew. Math. \textbf{53} (1857), 173--175.
\bibitem{Smyth} C.~J. Smyth, \emph{On the product of the conjugates outside the unit circle of an
algebraic integer}, Bull. London Math. Soc. \textbf{3} (1971), 169--175.
\bibitem{Dobrowolski} E.~Dobrowolski, \emph{On a question of Lehmer and the number of irreducible
factors of a polynomial}, Acta Arith. \textbf{34} (1979), 391--401.
\bibitem{Salem} R.~Salem, \emph{Power series with integral coefficients}, Duke Math. J. \textbf{12}
(1945), 153--172.
\bibitem{corpus} Echo--Squirrel Research, \emph{The Emission--Gap Theorem}, \emph{Lehmer's Box},
\emph{The Occupant of the Salem Slot}, and the substrate paper \texttt{vector\_substrate};
\texttt{math-research-pipelines} (2026). Companion stack to the present paper.
\end{thebibliography}

\smallskip
\noindent\footnotesize\textit{Provenance caveat.} These references are standard attributions
supplied from memory; volume, page, and year details have \emph{not} been machine-verified and
should be checked against the primary sources before external use. The Schur--Siegel--Smyth
attributions (Siegel~1945; Smyth~1984) name a recognized open problem and are cited for
orientation only.

\end{document}
